% GENERATED FILE. Edit canonical lessons, then run npm run generate:books.
% canonical-source-hash: fnv1a64:677208eefd4eb882
% canonical-lessons: IT-C484-radiografia, IT-C484-gesso, IT-C484-stampella, IT-C484-ingessare, IT-C484-zoppicare

\chapter{Una gamba ingessata}
\label{ch:it-una-gamba-ingessata}
\hlchaptermodality{484}

\begin{chapteropening}
\textbf{By the end of this chapter:} \emph{I can describe a fracture: the X-ray, the cast and the crutches.}

Complete the last lesson of chapter 484: I can describe a fracture: the X-ray, the cast and the crutches.
\end{chapteropening}

\section[la radiografia]{la radiografia — X-ray}

\label{lesson:IT-C484-radiografia}



\begin{quote}
Before the new one: say the Italian for to trip, then the Italian for to cut oneself.
\end{quote}

\subsection*{You'll want to know: la radiografia}
\textbf{la radiografia} — \textquotedblleft{}X-ray\textquotedblright{}.

\begin{quote}
Il medico mi ha fatto una radiografia al braccio.
\end{quote}

\begin{quote}
\emph{The doctor X-rayed my arm.}
\end{quote}

\emph{Fare una radiografia} — to have an X-ray. In speech often \emph{le lastre}.

\begin{cousinweb}
\emph{radiografia} is a modern learned word built from Latin \emph{radius} \textquotedblleft{}a ray\textquotedblright{} (Lewis \& Short) and Greek \emph{-graphía} \textquotedblleft{}writing\textquotedblright{} (EtymDB).
\end{cousinweb}

\begin{grammarlens}[title={What you've built}]
One more word for telling, explaining and describing at B1.
\end{grammarlens}

\subsection*{Your turn}
\begin{itemize}
\raggedright
  \item \emph{Say it:} \emph{la radiografia}
  \item \emph{Say it:} \emph{la radiografia}, once more
  \item \emph{Say it:} a sentence of your own with \emph{la radiografia}
  \item \emph{Recall:} say the Italian for to trip, then the Italian for to cut oneself, then say \emph{la radiografia} again
\end{itemize}

\begin{tcolorbox}[breakable,colback=teal!4,colframe=teal!35!black,title={Before you move on}]
What does la radiografia mean? (\textquotedblleft{}X-ray\textquotedblright{}.) Say it once more.
\end{tcolorbox}
\section[il gesso]{il gesso — plaster cast; chalk}

\label{lesson:IT-C484-gesso}



\begin{quote}
Before the new one: say the Italian for to cut oneself, then the Italian for X-ray.
\end{quote}

\subsection*{You'll want to know: il gesso}
\textbf{il gesso} — \textquotedblleft{}plaster cast; chalk\textquotedblright{}.

\begin{quote}
Porto il gesso al braccio per un mese.
\end{quote}

\begin{quote}
\emph{I am wearing a cast on my arm for a month.}
\end{quote}

At school, \emph{il gesso} is the chalk for the blackboard.

\begin{cousinweb}
\emph{gesso} is inherited from Latin \emph{gypsum} \textquotedblleft{}chalk, plaster\textquotedblright{} (EtymDB), from Greek \emph{gýpsos}.
\end{cousinweb}

\begin{grammarlens}[title={What you've built}]
One more word for telling, explaining and describing at B1.
\end{grammarlens}

\subsection*{Your turn}
\begin{itemize}
\raggedright
  \item \emph{Say it:} \emph{il gesso}
  \item \emph{Say it:} \emph{il gesso}, once more
  \item \emph{Say it:} a sentence of your own with \emph{il gesso}
  \item \emph{Recall:} say the Italian for to cut oneself, then the Italian for X-ray, then say \emph{il gesso} again
\end{itemize}

\begin{tcolorbox}[breakable,colback=teal!4,colframe=teal!35!black,title={Before you move on}]
What does il gesso mean? (\textquotedblleft{}Plaster cast; chalk\textquotedblright{}.) Say it once more.
\end{tcolorbox}
\section[la stampella]{la stampella — crutch}

\label{lesson:IT-C484-stampella}



\begin{quote}
Before the new one: say the Italian for X-ray, then the Italian for plaster cast.
\end{quote}

\subsection*{You'll want to know: la stampella}
\textbf{la stampella} — \textquotedblleft{}crutch\textquotedblright{}.

\begin{quote}
Cammino con le stampelle da due settimane.
\end{quote}

\begin{quote}
\emph{I have been walking on crutches for two weeks.}
\end{quote}

In some regions it is also a clothes hanger; elsewhere that is \emph{la gruccia}.

\begin{cousinweb}
The origin of \emph{stampella} is not settled.
\end{cousinweb}

\begin{grammarlens}[title={What you've built}]
One more word for telling, explaining and describing at B1.
\end{grammarlens}

\subsection*{Your turn}
\begin{itemize}
\raggedright
  \item \emph{Say it:} \emph{la stampella}
  \item \emph{Say it:} \emph{la stampella}, once more
  \item \emph{Say it:} a sentence of your own with \emph{la stampella}
  \item \emph{Recall:} say the Italian for X-ray, then the Italian for plaster cast, then say \emph{la stampella} again
\end{itemize}

\begin{tcolorbox}[breakable,colback=teal!4,colframe=teal!35!black,title={Before you move on}]
What does la stampella mean? (\textquotedblleft{}Crutch\textquotedblright{}.) Say it once more.
\end{tcolorbox}
\section[ingessare]{ingessare — to put in a plaster cast}

\label{lesson:IT-C484-ingessare}



\begin{quote}
Before the new one: say the Italian for plaster cast, then the Italian for crutch.
\end{quote}

\subsection*{You'll want to know: ingessare}
\textbf{ingessare} — \textquotedblleft{}to put in a plaster cast\textquotedblright{}.

\begin{quote}
In ospedale mi hanno ingessato la gamba.
\end{quote}

\begin{quote}
\emph{At the hospital they put my leg in a cast.}
\end{quote}

\emph{Ingessato} — in a cast: \emph{ho il braccio ingessato}.

\begin{cousinweb}
\emph{ingessare} is built on \emph{gesso} \textquotedblleft{}plaster\textquotedblright{}, inherited from Latin \emph{gypsum}, from Greek \emph{gýpsos} (EtymDB).
\end{cousinweb}

\begin{grammarlens}[title={What you've built}]
One more word for telling, explaining and describing at B1.
\end{grammarlens}

\subsection*{Your turn}
\begin{itemize}
\raggedright
  \item \emph{Say it:} \emph{ingessare}
  \item \emph{Say it:} \emph{ingessare}, once more
  \item \emph{Say it:} a sentence of your own with \emph{ingessare}
  \item \emph{Recall:} say the Italian for plaster cast, then the Italian for crutch, then say \emph{ingessare} again
\end{itemize}

\begin{tcolorbox}[breakable,colback=teal!4,colframe=teal!35!black,title={Before you move on}]
What does ingessare mean? (\textquotedblleft{}To put in a plaster cast\textquotedblright{}.) Say it once more.
\end{tcolorbox}
\section[zoppicare]{zoppicare — to limp}

\label{lesson:IT-C484-zoppicare}



\begin{quote}
Before the new one: say the Italian for crutch, then the Italian for to put in a plaster cast.
\end{quote}

\subsection*{You'll want to know: zoppicare}
\textbf{zoppicare} — \textquotedblleft{}to limp\textquotedblright{}.

\begin{quote}
Dopo la partita zoppicava un po'.
\end{quote}

\begin{quote}
\emph{After the match he was limping a little.}
\end{quote}

Built on \emph{zoppo} \textquotedblleft{}lame\textquotedblright{}.

\begin{cousinweb}
\emph{zoppicare} is built on \emph{zoppo} \textquotedblleft{}lame\textquotedblright{} (EtymDB). The origin of \emph{zoppo} is not settled.
\end{cousinweb}

\begin{grammarlens}[title={What you've built}]
One more word for telling, explaining and describing at B1.
\end{grammarlens}

\subsection*{Your turn}
\begin{itemize}
\raggedright
  \item \emph{Say it:} \emph{zoppicare}
  \item \emph{Say it:} \emph{zoppicare}, once more
  \item \emph{Say it:} a sentence of your own with \emph{zoppicare}
  \item \emph{Recall:} say the Italian for crutch, then the Italian for to put in a plaster cast, then say \emph{zoppicare} again
\end{itemize}

\begin{tcolorbox}[breakable,colback=teal!4,colframe=teal!35!black,title={Before you move on}]
What does zoppicare mean? (\textquotedblleft{}To limp\textquotedblright{}.) Say it once more.
\end{tcolorbox}
